%%%%%%%%%%%%
%
% $Autor: Wings $
% $Datum: 2019-03-05 08:03:15Z $
% $Pfad: SetUp.tex $
% $Version: 4250 $
% !TeX spellcheck = en_GB/de_DE
% !TeX encoding = utf8
% !TeX root = manual 
% !TeX TXS-program:bibliography = txs:///biber
%
%%%%%%%%%%%%

\chapter{Aufbau- und Inbetriebnahmeplan}

\subsection{Schritt 1: Standortwahl und Vorbereitung}

\begin{itemize}
	\item \textbf{Untergrund:} Wählen Sie eine ebene, stabile und trockene Arbeitsfläche (Tisch oder Werkbank).
	
	\item \textbf{Platzbedarf:} Achten Sie darauf, dass an beiden Enden des Bandes genügend Platz für die Waren oder Auffangbehälter vorhanden ist.
	
	\item \textbf{Umgebung:} Stellen Sie sicher, dass keine losen Gegenstände oder Kabel in die Nähe der Laufrollen gelangen können.
\end{itemize}

\subsection{Schritt 2: Mechanischer Aufbau}

\begin{itemize}
	\item \textbf{Ausklappen:} Klappen Sie die Stützbeine nacheinander vollständig aus.
	
	\item \textbf{Ausrichtung:} Kontrollieren Sie mit einer Wasserwaage, ob das Band gerade steht. Ein schiefes Band führt dazu, dass das Transportgut zur Seite rutscht.
\end{itemize}

\subsection{Schritt 3: Elektrischer Anschluss}

\begin{itemize}
	\item \textbf{Sichtprüfung:} Überprüfen Sie das Gehäuse der Steuereinheit und das Netzteil auf sichtbare Schäden.
	
	\item \textbf{Verbindung:} Stecken Sie den DC-Stecker des Netzteils in die Buchse an der Steuereinheit des Förderbandes.
	
	\item \textbf{Netzanschluss:} Stecken Sie das Netzteil in eine gut erreichbare Steckdose (100--240 V).
	
	\item \textbf{Hinweis:} Das Gerät sollte sofort einsatzbereit sein, aber noch nicht anlaufen.
\end{itemize}

\subsection{Schritt 4: Justierung und Transportband (Vor dem ersten Start)}

\begin{itemize}
	\item \textbf{Mittigkeit:} Prüfen Sie, ob das Transportband mittig auf den Rollen liegt.
	
	\item \textbf{Spannung:} Drücken Sie leicht mit dem Finger auf die Mitte des Transportbands. Es sollte stramm sitzen, aber eine minimale Flexibilität (ca. 0,5--1 cm Spiel) haben.
	
	\item \textbf{Fremdkörper:} Stellen Sie sicher, dass keine Werkzeuge oder Verpackungsreste auf dem Transportband liegen.
\end{itemize}

\subsection{Schritt 5: Inbetriebnahme und Testlauf}

\begin{itemize}
	\item \textbf{Modus wählen:}
	Wählen sie für die Funktionsprüfung mit dem Modusschalter den Modus L aus.  
	
	\item \textbf{Geschwindigkeit auf Minimum:} Drehen Sie den Drehregler auf die kleinste Stufe.
	
	\item \textbf{Starten:} Bitte legen sie einen faustgroßen Gegenstand auf die rechte Seite des Transportbandes, in die Lichtschranke. 
	
	\item \textbf{Funktionsprüfung:}
	\begin{itemize}
		\item Läuft das Transportband ruhig und ohne schleifende Geräusche?
		\item Bleibt das Transportband in der Mitte oder wandert er zu einer Seite aus? (Falls ja: Justierschrauben an den Laufrollen vorsichtig anpassen).
	\end{itemize}
	
	\item \textbf{Geschwindigkeitstest:} Erhöhen Sie langsam die Geschwindigkeit, um die Stabilität der Stützbeine bei Vibrationen zu testen.
\end{itemize}

\subsection{Schritt 6: Regelbetrieb (Laufrichtng L oder R)}

\begin{itemize}
	\item \textbf{Aufsichtspflicht:} Bleiben Sie während des Betriebs immer in Reichweite der Hauptschalter.
	
	\item \textbf{Laufrichtung:} Über den Modusschalter [TeilNr03] lässt sich die Laufrichtung jederzeit zwischen [L] und [R] umschalten. Bei einem Richtungswechsel im Betrieb bremst das Band kontrolliert ab und beschleunigt anschließend in die Gegenrichtung. Da das System am Ende nicht automatisch stoppt, muss das Transportgut rechtzeitig manuell entnommen werden. Ergänzende Details hierzu finden Sie im [Kapitel 6.1 (Umkehr der Laufrichtung)] .
	
	\item \textbf{Not-Halt:} Im Falle einer Blockade drücken Sie sofort die Stopp-Taste. Greifen Sie niemals bei laufendem Motor in das Transportband.
\end{itemize}

\subsection{Schritt 7: Regelbetrieb (Modus A)}

\begin{itemize}
	\item \textbf{Aufsichtspflicht:} Bleiben Sie während des Betriebs immer in Reichweite der Hauptschalter.
	
	\item \textbf{Automatik-Manöver:} Im Automatikmodus befördert das Transportförderband Objekte nach der Erkennung selbstständig bis zur gegenüberliegenden Lichtschranke. Die Bandgeschwindigkeit lässt sich dabei variabel zwischen 1,8 und 7,6 cm/s einstellen und bleibt auch unter Last konstant. Integrierte Funktionen sorgen für einen automatischen Stopp bei Zeitüberschreitung sowie einen zeitgesteuerten System-Reset.Weitere technische Details und Bedienhinweise finden Sie im [Kapitel 6.1 (Automatikbetrieb)].
	
	\item \textbf{Not-Halt:} Im Falle einer Blockade drücken Sie sofort die Stopp-Taste. Greifen Sie niemals bei laufendem Motor in das Transportband.
\end{itemize}